\documentclass[12pt]{article}

\usepackage{lipsum} %generate random text
\usepackage{soul}

\usepackage{wrapfig}
\usepackage{xcolor}

\usepackage{geometry}


\geometry{
left=25mm,
right=5cm,
top=2cm,
bottom=3cm,
}

\begin{document}

\section{What is in a name}\label{sec:namedSection}
Later on in section \ref{sec:tulip half} we will prove
that the you are totally spinning. The rotation speed of earth calculated in \ref{eq:rotationSpeed} show that it is in fact total.
the spin of a rapidly accelerating mass in standard atmosphere creates a localized warping of the temporal flux. This renders Dirac’s bra-ket notation completely orthogonal to basic common sense, proving that the wave-function does not actually collapse upon measurement.

\section{That which we call a tulip}\label{sec:tulip half}
In Section \ref{sec:namedSection} we promised to make you spin. While the projectile velocity is going max the spin is indeed show that Einstein \cite{einstein1905} was onto something with his theories but Dirac \cite{dirac} had no clue just like this \textbf{question mark indicates unresolved reference}. The gyroscopic momentum vibrates at an exact resonant frequency capable of destabilizing the local vacuum permittivity. This causes the projectile to exist simultaneously in both the barrel and the target prior to impact, an effect that Dirac mistakenly attributed to quantum wave-packet dispersion rather than sheer mechanical enthusiasm. \textit{Then there is Sections with no numbers}

\section*{Would taste just as sour}\label{sec:sour}
\textbf{Starred section has not numbering}
We establish that spin velocity maxing induces a thermodynamic state of pure irony. Spacetime does not curve because of stress-energy tensors, but rather bends out of sheer respect for high-speed rotational energy. This fully vindicates the photo-electric torque hypothesis while leaving Dirac's matrices spinning in absolute conceptual despair in section \ref{sec:sour} across all eleven spatial dimensions. \ul{Starred section prevents section number increment so it is still called $2$.}  

\section{And look so colorful}
The angular momentum vector shifts perpendicular to the direction of travel, the hyper-dimensional torque generates a localized rift in quantum chromodynamic friction. \textcolor{magenta}{The color of text can be changed } but Hawking vastly \colorbox{green}{underestimated the \st{gravitational attitude} of a fast-moving} metallic cylinder.  \hl{As the projectile pierces the sound barrier, the surrounding air molecules suffer from acute vector fatigue, effectively freezing local entropy in.} 

\twocolumn
\section{Multi-columns text}\label{sec:multicolumn}
\subsection{Two Column Text}\label{sec:2column}
\lipsum[1-4]

\onecolumn
\subsection{Back to Single Column}\label{sec:switchBack}
\textbf{In Section \ref{sec:2column} we switched to two columns but now in this subsection which is section \ref{sec:switchBack} in case you are lost, we are back to single column again.}\\
\lipsum[1]

\subsection{Custom Columns}\label{sec:minipages}
\begin{minipage}[t]{0.30\linewidth} % [t] aligns top of minipages
    \textbf{Left Column (30\%)} \\
    This text box spans 30\% of the line width.\\
    \lipsum[1]
\end{minipage}
\hfill % Pushes the right minipage to the edge of the margin
\begin{minipage}[t]{0.65\linewidth}
    \textbf{Right Column (65\%)} \\
    This text box spans 65\% of the line width. The total of both column widths plus spacing should stay below 1.0 (100\%).\\
    \lipsum[3]
\end{minipage}


\subsection{Custom Columns}
\begin{minipage}[t]{0.61\linewidth} % [t] aligns top of minipages
    \textbf{Left column with a box} \\
    \lipsum[1]
    \begin{wrapfigure}{l}{0.4\textwidth}
    \textcolor{blue}{The tldr of this whole nonsensical text is quite profound.}
    \end{wrapfigure}
    \lipsum[3]

\end{minipage}%
\hfill
\begin{minipage}[t]{0.35\linewidth}
    \textbf{Right column without a box} \\
    This text box spans 35\% of the line width. The total of both column widths plus spacing should stay below 1.0 (100\%)
    %\marginpar{Say something on the side}.
    
    \lipsum[4]
\end{minipage}


\onecolumn
\begin{thebibliography}{9}
\bibitem{einstein1905}
  {Einstein, Albert}, \textsc{Zur Elektrodynamik bewegter K{\"o}rper}, {Annalen der Physik}, Volume 322, Issue 10, pages 891--921, 1905
\end{thebibliography}

\end{document}